\subsection{APPLICATION TO THE EXTENSION OF REAL-VALUED FUNCTIONS}

Since \(\mathbf{R}\) can be considered as a subset of \(\tilde{\mathbf{R}}\) , every mapping of a set \(\mathbf{E}\) into \(\mathbf{R}\) (i.e.~every real-valued function on \(\mathbf{E}\) ) can be considered as a mapping of \(\mathbf{E}\) into \(\tilde{\mathbf{R}}\) ; in particular, if \(\mathbf{E}\) is a subset of a topological space \(\mathbf{F}\) and \(f\) is a mapping of \(\mathbf{E}\) into \(\mathbf{R}\) , it may happen that at some points of the closure \(\overline{\mathbf{E}}\) of \(\mathbf{E}\) , \(f(x)\) tends to the limit \(\infty\) as \(x\) tends to one of these points while remaining in \(\mathbf{E}\) ; we may then extend the function \(f\) by continuity by assigning it the value \(\infty\) at these points (Chapter I, § 8, no. 5, Theorem 1).

Consider in particular the case where E is a subset of \(R^{n}\) , the space \(R^{n}\) itself being considered as embedded in projective space \(P_{n}\) ; if we suppose that the hyperplane at infinity is \(x_{0}=0\) , a real-valued function f defined on E may be identified with the mapping

\[
(x _ {0}, x _ {1}, x _ {2}, \ldots , x _ {n}) \rightarrow f \left(\frac {x _ {1}}{x _ {0}}, \frac {x _ {2}}{x _ {0}}, \ldots , \frac {x _ {n}}{x _ {0}}\right)
\]

of E into \(\tilde{R}\) ; from what has been said in the previous paragraph it may be possible to extend this function, not only to some points of \(R^{n}\) in the closure of E, but also to some of the ``points at infinity'' of \(P_{n}\) in the closure of E.

Let us show that we obtain in this way, for example, the extension by continuity to the whole of \(\tilde{\mathbf{R}}\) of a rational function of a real variable, already defined in algebra. Let us identify \(\tilde{\mathbf{R}}\) and \(\mathbf{P}_1\) , every real number \(x \in \mathbb{R}\) being identified with the point whose homogeneous coordinates are (1, x), and the point \(\infty\) with the point whose homogeneous coordinates are (0,1). Let \(u(x)/v(x)\) be a rational function, \(u\) and \(v\) being two coprime polynomials of degrees \(m\) and \(n\) respectively; if we suppose for example that \(m \leqslant n\) , and if we put \(u_1(x,y) = x^n u(y/x)\) , \(v_1(x,y) = x^n v(y/x)\) , the rational function \(u/v\) may be considered as the restriction, to the set of real numbers \(x\) such that \(v(x) \neq 0\) , of the mapping \((x,y) \to (v_1(x,y), u_1(x,y))\) . In other words, we extend \(u/v\) by continuity, by giving it the value \(\infty\) at those points \(x \in \mathbb{R}\) where \(v(x) = 0\) , and by giving it at the point \(\infty\) the value \(0\) if \(m < n\) , the value \(\infty\) if \(m > n\) , and the value of the ratio of the leading coefficients if \(m = n\) .

In particular, the function \(1 / x\) can be extended to the point \(o\) by taking the value \(\infty\) there, to \(\infty\) by taking the value \(o\) there; this extended function is evidently a homeomorphism of \(\tilde{\mathbf{R}}\) onto itself. The same is true of the homographic function \((ax + b) / (cx + d)\) when \(ad - bc \neq o\) .

Likewise, if \(n\) is an integer \(> 0\) , the function \(x^n\) extends to the point \(\infty\) by taking the value \(\infty\) there.

On the other hand, it is in general impossible to extend by continuity a rational function of two real variables either to the space \(\mathbf{P}_1 \times \mathbf{P}_1\) or to the space \(\mathbf{P}_2\) (cf.~Exercise 5).

